\documentclass[11pt,a4paper]{article}
\usepackage[margin=1in]{geometry}
\usepackage{amsmath,amssymb,amsthm}
\usepackage{algorithm}
\usepackage{algpseudocode}
\usepackage{booktabs}
\usepackage{hyperref}

\theoremstyle{definition}
\newtheorem{definition}{\textbf{Definition}}[section]
\newtheorem{theorem}{\textbf{Theorem}}[section]
\newtheorem{lemma}[theorem]{\textbf{Lemma}}
\newtheorem{remark}{\textbf{Remark}}[section]

\title{\textbf{ALGO-NN-271: Integrated gradients and gradient saliency}}
\author{\textbf{Team Scaibu} \\ \small Email: \texttt{jaydeep.9664920749@gmail.com} \\ \small \textbf{Architecture and Core Engine Team}}
\date{\textbf{October 2026}}

\begin{document}
\maketitle

\section{Definitions and Cognitive Mental Models}

\subsection{Formal Mathematical Definition}
\textbf{Definition (Algorithm \#271 --- Integrated gradients and gradient saliency):}
Attributing a prediction to input features by accumulating gradients along a path from a baseline input to the actual input.

Formally, given continuous operational inputs $\mathbf{x} = [x_1, x_2, \dots, x_d]^T \in \mathbb{R}^d$ and parameter manifold $\Theta$, the transformation $\mathcal{T}_{\text{algo}}(\mathbf{x}; \Theta)$ satisfies:
{\boldmath
\begin{equation}
\mathbf{y} = \mathcal{T}_{\text{algo}}(\mathbf{x}; \Theta) = \arg\min_{\mathbf{y} \in \mathcal{Y}} \mathcal{L}(\mathbf{x}, \mathbf{y}) + \lambda \mathcal{R}(\mathbf{y})
\end{equation}
}
where $\mathcal{L}$ quantifies reconstruction fidelity and $\mathcal{R}$ enforces algorithmic structural constraints.

\subsection{Expanded Technical Specification}
\begin{itemize}
\item \textbf{**Plain gradient saliency:** |}: ∂output / ∂input|. Simple, but noisy, and suffers from saturation (gradients near zero even for important features).
\item \textbf{**Integrated gradients:** choo}: se a baseline x′ (for example a black image or zero embeddings). Attribution for feature i = (x\_i − x′\_i) × the average of ∂output/∂x\_i along the straight path from x′ to x (approximated with 20–300 steps).
\item \textbf{**Completeness:** attributions}:  sum to output(x) − output(x′).
\item \textbf{**SmoothGrad:** average gradie}: nts over noisy copies of the input to reduce noise.
\end{itemize}

\subsection{Computational Complexity and Cost Model}
\begin{equation}
\mathcal{O}(T \cdot d \cdot K) \quad \text{computational operations}
\end{equation}
\textbf{Operational Resource Cost:} Tens to hundreds of backward passes per explanation.

\section{Theoretical Proof and Algorithmic Execution}

\begin{algorithm}[H]
\caption{Integrated gradients and gradient saliency Operational Pipeline}
\begin{algorithmic}[1]
\Require Input tensor $\mathbf{x} \in \mathbb{R}^d$, hyperparameter set $\mathcal{H}$
\Ensure Optimal processed tensor $\mathbf{y} \in \mathbb{R}^d$, diagnostic metric record $\mathcal{M}$
\State Initialize operational parameters $\Theta_0 \leftarrow \text{Extract}(\mathcal{H})$
\State $\mu \leftarrow \frac{1}{d} \sum_{i=1}^d x_i$, \quad $\sigma^2 \leftarrow \frac{1}{d} \sum_{i=1}^d (x_i - \mu)^2$
\For{each channel or coordinate $j \in \{1, \dots, d\}$}
    \State Evaluate local transformation invariant $\hat{x}_j \leftarrow \phi(x_j, \mu, \sigma, \mathcal{H})$
    \State Compute bounded output $y_j \leftarrow \text{Clip}(\hat{x}_j, \text{bounds}(\mathcal{H}))$
\EndFor
\State $\mathcal{M} \leftarrow \text{ComputeDiagnostics}(\mathbf{x}, \mathbf{y})$
\State \Return $\mathbf{y}, \mathcal{M}$
\end{algorithmic}
\end{algorithm}

\begin{theorem}[Convergence and Bounded Error]
Under Lipschitz continuous parameter mapping with constant $L < \infty$, the output distortion satisfies:
\begin{equation}
\| \mathbf{y} - \mathbf{x}^* \|_2 \le \frac{L}{1 - \gamma} \epsilon_{\text{quant}}
\end{equation}
guaranteeing numeric stability across execution cycles.
\end{theorem}

\begin{proof}
By contraction mapping over the bounded metric space $(\mathcal{X}, \|\cdot\|_2)$, successive projections satisfy $\|\mathcal{T}(\mathbf{x}) - \mathcal{T}(\mathbf{x}')\|_2 \le \gamma \|\mathbf{x} - \mathbf{x}'\|_2$ with $\gamma < 1$. Unrolling the recurrence confirms exponential decay of truncation error.
\end{proof}

\section{Production Systems Architecture and Agent Policy}
\begin{itemize}
\item \textbf{Role Integration}: Autonomous agents consume this operator during automated validation, inference scaling, and policy verification.
\item \textbf{Guardrails and Rules}: Strict dynamic bounds prevent numerical underflow and activation divergence in distributed serving fleets.
\end{itemize}

\end{document}
